% ============================================================
%  Response to Reviewers, PPAM 2026, submission 25
%  Every number quoted here is reproducible from
%  artifacts/tidy/runs.csv via this folder's analysis/analyse.py.
% ============================================================
\documentclass[11pt,a4paper]{article}

\usepackage[T1]{fontenc}
\usepackage[utf8]{inputenc}
\usepackage[a4paper,margin=2.4cm]{geometry}
\usepackage{microtype}
\usepackage{siunitx}
\usepackage{booktabs}
\usepackage{enumitem}
\usepackage{xcolor}
\usepackage[colorlinks=true,allcolors=blue]{hyperref}

\sisetup{detect-weight=true, detect-family=true}
\DeclareSIUnit\mebibyte{MiB}
\emergencystretch=3em

\newcommand{\tmma}{\texttt{torch.cuda.\allowbreak max\_memory\_\allowbreak allocated}}
\newcommand{\FF}{FF}
\newcommand{\MF}{MF}
\newcommand{\BP}{BP}
\newcommand{\CaFo}{CaFo}
\newcommand{\CaFoRand}{CaFo-Rand}
\newcommand{\CaFoDFA}{CaFo-DFA}
\newcommand{\BPDS}{BP-DS}
\newcommand{\MFJ}{MF-Joint}

\newenvironment{comment}%
  {\begin{quote}\itshape\small}%
  {\end{quote}}
\newcommand{\where}[1]{\par\smallskip\noindent\textbf{Where:} #1\par}
\setlist[description]{leftmargin=0pt, labelindent=0pt, style=nextline}

\title{\vspace{-1.5cm}Response to Reviewers\\[2pt]
  \large Energy-Efficient Deep Learning Without Backpropagation:\\
  A~Rigorous Hardware-Validated Benchmarking Study of Forward-Only Algorithms\\[4pt]
  \normalsize PPAM 2026, Submission 25 (CP025)}
\author{}
\date{}

\begin{document}
\maketitle
\vspace{-1.2cm}

% ============================================================
\section*{Summary}
% ============================================================

We thank all five reviewers. Two of the comments did more than improve the
paper: they produced the experiment that the revised version is built on, and
they changed our headline result. We start with that one.

\paragraph{R3c and R5d asked the decisive question, and the answer is in
\MF{}'s favour.}
Both reviewers observed that \MF{} places a cross-entropy loss at every hidden
layer, so it behaves like a deep-supervision network, and that without a
matched deep-supervision baseline we could not tell whether \MF{}'s accuracy
comes from forward-only training or merely from localised objectives. That
control now exists. \BPDS{} is \BP{} carrying \MF{}'s auxiliary classifiers
with nothing detached, and on the two CIFAR \num{3}$\times$\num{2000} MLPs it
is \emph{worse} than plain \BP{}, by \num{1.65}~pp on CIFAR-10 and
\num{1.06}~pp on CIFAR-100. \MF{} beats that control by \num{5.56}~pp
$[\num{+5.37},\num{+5.75}]$ and \num{6.10}~pp $[\num{+5.88},\num{+6.32}]$ over
\num{20} paired seeds. The localised objective does not explain the gain;
adding it to \BP{} lowers \BP{}'s accuracy. The evidence therefore isolates
gradient detachment/locality as the dominant tested factor; we did not test
every possible mechanism. On the two MNIST-scale
\num{2}$\times$\num{1000} MLPs the same contrast is a tie, and we report that
half too. The full four-rung ladder is Section~5.3,
Table~3 and Figure~2 of the revised paper.

\paragraph{R2d asked whether \BP{} simply has better kernels, and the
telemetry answers it.}
We tested the objection rather than arguing it. On the four ladder
configurations no method exceeds \SI{20.3}{\percent} mean GPU utilisation and
\BP{} itself averages \SI{11.5}{\percent}; across the whole study, including
the wider MLPs used for \FF{}, no run exceeds
\SI{38.0}{\percent}. Every arm is launch-bound, not arithmetic-bound, so
\BP{} does not win by saturating the A100 with more efficient arithmetic. We also built
the obvious missing optimisation, an activation cache that removes \MF{}'s
$O(L^2)$ recomputation, verified it bit-exact against the recomputing path
(maximum absolute difference in loss and weights exactly \num{0} over
\num{1200} and \num{3600} optimiser steps), and measured that it saves no time
while costing \SI{374}{\percent} and \SI{404}{\percent} more per-process
allocator memory. These measurements do not entirely exclude an advantage
from implementation maturity, but the difference they do measure between the
methods lies in convergence rate rather than in the cost of a pass.

\paragraph{The cost result has two readings, and this version states both.}
Trained to its own convergence, \MF{} costs \SIrange{3.2}{4.0}{\times}
\BP{}'s energy. The reason it costs more end to end is that it needs
\SIrange{3.6}{4.4}{\times} the wall-clock time, not that each pass is
expensive. Given the number of data passes \BP{} actually spent, \MF{} is the
\emph{cheaper} algorithm: \SI{19.9}{\percent} less GPU energy on CIFAR-10 and
\SI{12.5}{\percent} less on CIFAR-100, \SI{27.5}{\percent} and
\SI{24.6}{\percent} less per-process allocator memory, at wall-clock time within
\SI{3}{\percent} of \BP{}'s. The revised paper leads with that
matched data-pass protocol result (Section~5.1, Table~2), keeps the
own-convergence result in full (Section~5.4, Table~4), and names the
pass multiplier as the open problem (Section~6).

\paragraph{What is new in this version, concretely.}
\begin{itemize}[topsep=4pt,itemsep=0pt,parsep=0pt,leftmargin=1.4em]
  \item \textbf{A matched data-pass protocol}, in which \MF{} receives
        $\lceil \bar{E}_{\BP{}}/S \rceil$ epochs per stage, so rounding can
        only give \MF{} more data passes than \BP{} had. The protocol
        equalises data passes, not floating-point operations or optimiser
        updates per layer, and is named accordingly. Section~4, Table~1.
  \item \textbf{A pre-registered four-rung ablation ladder} from \BP{} to
        \MF{}, including the \BPDS{} control of R3c and R5d, over
        \numrange{20}{122} seeds per rung. Section~5.3.
  \item \textbf{A new figure showing that the advantage is larger in the
        larger configurations tested}, from no memory advantage at \num{1.8}~M weights to
        \SI{27.5}{\percent} at \num{14.2}~M. Figure~1.
  \item \textbf{Seed counts of \numrange{20}{122} per \MF{} contrast} in place of
        three, with every parity claim resolved by the two one-sided tests
        (TOST) procedure against a pre-registered $\pm\num{0.25}$~pp margin,
        now cited and explained in Section~4. R4b, R5c.
  \item \textbf{Every figure regenerated} from the measurement archive by a
        script in the repository. The illegible dashboard exports of R1b are
        gone. R1b, R3b.
\end{itemize}

\noindent
Every measured number in the revised paper is produced from a single tidy table
of \num{19944} measurements by \texttt{analysis/analyse.py}, which is in the
repository alongside the figure script; only the realised epoch counts are
transcribed, from the run configurations. The manuscript is \textbf{15 pages},
the PPAM limit.

\bigskip
\noindent\textbf{Revised structure.}
\S1~Introduction \textbar{} \S2~Background and Related Work \textbar{}
\S3~Mono-Forward and the Cost Hypothesis \textbar{}
\S4~Measuring Cost Fairly: Two Protocols \textbar{}
\S5~Results (5.1 the matched data-pass protocol, 5.2 scale, 5.3 the ladder,
5.4 own convergence, 5.5 quantifying two objections, 5.6 \FF{} and \CaFo{})
\textbar{} \S6~Discussion and Limitations \textbar{} \S7~Conclusion.
Table~1 = budgets; Table~2 = matched data-pass protocol; Table~3 = the ladder;
Table~4 = own convergence. Figure~1 = saving against model scale;
Figure~2 = the ladder.

% ============================================================
\section*{Reviewer 1}
% ============================================================

\subsection*{R1a: ``Reads like a very condensed technical report \dots\
detrimental to readability for the average PPAM reader.''}

\begin{comment}
Reads like a very condensed technical report, which is detrimental to
readability for the average PPAM reader.
\end{comment}

\noindent
Accepted, and the paper has been rewritten rather than edited. The accepted
version was organised as one subsection per algorithm, which forced the reader
to assemble the argument themselves. The revision is organised around one
claim and its evidence: Section~3 states a mechanism and why it should be
cheaper, Section~4 states the two protocols and why two are needed, Section~5
tests the claim, and Section~6 says where it fails. Three of the six results subsections (5.1, 5.2 and 5.4) carry a verdict in their titles, and the three emphasised sentences of
Section~5.1 (\emph{the memory figures are exact}, \emph{the energy saving is a
power saving}, \emph{on the most demanding task the saving costs no
accuracy}) give a reader
the result without the table.

Two of the narrowing devices are deliberate. We dropped the
algorithm-by-algorithm survey structure and compressed \FF{} and \CaFo{} into
a single subsection (5.6), because the study's defensible contribution is
about \MF{} and the other two are context for it. We also state the scope of
every headline number in the sentence that carries it, rather than deferring
all qualification to a limitations section.

\where{Whole manuscript; Sections~3 to 6 are new text.}

\subsection*{R1b: ``The plots in Fig.~3 and Fig.~4 are completely illegible.''}

\begin{comment}
The plots in Fig.~3 and Fig.~4 are completely illegible.
\end{comment}

\noindent
Accepted. Those figures were raw dashboard exports at a size no print process
could rescue. They have been deleted, not enlarged.

The revision's figures are generated by
\texttt{analysis/figures.py} from the same tidy table as the tables, at the
text width of the \texttt{llncs} block, with every label in Figures~1 and~2 at \num{8}~pt or larger and no element carrying information by colour alone. Series are
separated by fill, hatch and line style together, so the figures survive a
monochrome print, which we checked by rendering them in greyscale. Figure~1 is
a four-group bar chart and Figure~2 is two line panels over four ladder rungs;
neither has more than four series. A third figure showing both protocols side
by side is generated by the same script and is in the repository, left out of
the manuscript only for space, since Table~4 carries the same numbers.

\where{Figures~1 and 2 and their captions; the supplementary third figure;
\texttt{analysis/figures.py}.}

% ============================================================
\section*{Reviewer 2}
% ============================================================

\subsection*{R2a: ``Lack of details and analysis in the main article
(specifically Section~6) made the paper seem preliminary.''}

\begin{comment}
Lack of details and analysis in the main article (specifically Section~6) made
the paper seem preliminary.
\end{comment}

\noindent
Accepted. The discussion is now built from measurements rather than from
interpretation, and it names what it cannot support.

It is built in two steps, followed by the implications; the robustness of
the cost result (the matched data-pass protocol gives \MF{} at least as many
data passes as \BP{} used, and board power is a rate that no stopping rule can
influence) is stated where it applies: the budget rule in Sections~1 and~4,
the power argument with the result itself in Section~5.1. \emph{Why gradient
detachment might improve accuracy} advances the regularisation reading, ties it
to the ladder result, and labels it a hypothesis that four configurations
cannot establish. \emph{The open problem: convergence rate, not per-unit cost}
identifies convergence
rate as the binding constraint, shows that the two obvious remedies fail
(caching is bit-exact and saves no time; the runs are launch-bound so faster
kernels would be an unlikely remedy, although implementation maturity is not
entirely excluded), and bounds the result on three further axes: scope,
accuracy under the matched data-pass protocol, and hardware. What was pre-registered and
what was not is stated in Section~4.

We also state plainly where the measurements and the narrative pulled against
each other. Under the matched data-pass protocol \MF{}'s energy is lower on all four
configurations by point estimate, a difference resolved on three, but it is as
accurate as \BP{} on only one of the four; we say so in Section~5.1 and
repeat it in the limitations, because it bounds the claim.

\where{Section~6, all paragraphs.}

\subsection*{R2b: ``The details of the different ML architectures are
important \dots\ these details are not discussed much in the paper.''}

\begin{comment}
The details of the different ML architectures are important, as they can
affect the results; these details are not discussed much in the paper.
\end{comment}

\noindent
Accepted, and the revision goes further than describing them: it makes
architecture an axis of the result.

Section~3 gives \MF{}'s construction explicitly, including the fact that an
MLP with $L$ hidden layers trains $L+1$ stages, one per projection matrix
including $\mathbf{M}_0$ on the raw input, which is what sets the epoch budget in
Table~1. Section~2 now introduces the notation and briefly describes \FF{}, \CaFo{}
and \MF{} before they are used. Section~4 states initialisation (Kaiming
uniform for all weights and projections, now with references for it and for
both optimisers), the optimiser per algorithm (AdamW for \BP{} and \BPDS{}, Adam
for \MFJ{} and \MF{} stages), the activation (ReLU) and batch size (\num{128}), the tuning budget and search width per algorithm, and the augmentation. It records that
every algorithm is built on the same \texttt{nn.Linear} primitives in one
codebase, differing only in the loss construction and in where
\texttt{detach()} is called. The remaining per-algorithm implementation
constants, the \CaFo{} block ($3{\times}3$ convolution, ReLU, $2{\times}2$
max-pooling, batch normalisation, \num{32}/\num{128}/\num{512} channels, three
block predictors summed) and the \FF{} goodness threshold $\theta=\num{2.0}$
with peer normalisation, are in the released configurations, since they bear
on \FF{} and \CaFo{} rather than on the claim.

Section~5.2 then uses the architectures. The four configurations span
\num{1.79}~M to \num{14.34}~M trainable weights, and every cost axis improves
across that range: energy saving from \SI{2.2}{\percent} and
\SI{4.3}{\percent} to \SI{12.5}{\percent} and \SI{19.9}{\percent}, memory from
a \SI{0.3}{\percent} deficit to a \SI{24.6}{\percent} and \SI{27.5}{\percent}
saving, so the advantage is larger in the larger configurations tested. We
are explicit that width is confounded with depth and with task across these
four, so the figure cannot separate width from scale, and Section~6 states that
the headline percentages hold on the \num{3}$\times$\num{2000} MLPs only.

\where{Section~3; Section~4, ``What else is held fixed''; Section~5.2 and
Figure~1.}

\subsection*{R2c: ``The paper claimed that the BP baselines were identical
but algorithm-specific components were removed. The removal of these
components may be the reason why certain architectures were better one way or
another.''}

\begin{comment}
The paper claimed that the BP baselines were identical but algorithm-specific
components were removed. The removal of these components may be the reason why
certain architectures were better one way or another.
\end{comment}

\noindent
This is correct, it was a real hole, and the ladder was built to close it. We
did not argue that the removal was harmless; we put the components back one at
a time and measured what each is worth.

\begin{center}\small
\begin{tabular}{llcll}
  \toprule
  Rung & Readout & Auxiliary losses & Gradients & Optimiser \\
  \midrule
  \BP{}    & output layer & no  & global & AdamW \\
  \BPDS{}  & output layer & yes & global & AdamW \\
  \MFJ{}   & $\mathbf{M}_L$ & yes & global & Adam \\
  \MF{}    & $\mathbf{M}_L$ & yes & \emph{local} & Adam \\
  \bottomrule
\end{tabular}
\end{center}

\noindent
Each adjacent pair differs in one structural property. We should be precise
about what that does and does not license: each rung is tuned independently
under the same \num{50}-trial budget, so every step also moves the learning
rate and weight decay, and the \BPDS{}$\rightarrow$\MFJ{} step additionally
crosses AdamW to Adam. That step is therefore not attributable to the readout
alone, and we no longer claim it is. The step our mechanism claim rests on,
\MFJ{}$\rightarrow$\MF{}, holds the optimiser family fixed at Adam. Accuracy
against \BP{}, in percentage points, as a difference of arm means over
\numrange{20}{122} seeds:

\begin{center}\small
\begin{tabular}{lrrrr}
  \toprule
  & MNIST 2$\times$1000 & Fashion 2$\times$1000 & CIFAR-10 3$\times$2000 & CIFAR-100 3$\times$2000 \\
  \midrule
  \BPDS{} & $+0.07$ & $+0.04$ & $-1.65$ & $-1.06$ \\
  \MFJ{}  & $+0.14$ & $-0.09$ & $-0.19$ & $+1.86$ \\
  \MF{}   & $+0.19$ & $+0.01$ & $+3.90$ & $+5.04$ \\
  \bottomrule
\end{tabular}
\end{center}

\noindent
Adding the auxiliary classifiers to \BP{} does not explain \MF{}'s advantage;
on the two CIFAR tasks it costs \BP{} accuracy. The readout step recovers most
or all of the loss, $+\num{1.46}$ and $+\num{2.92}$~pp, but it is the step that also
crosses optimisers, so we do not attribute it to the readout. The largest
single step, and the only one we attribute to a mechanism, is the detach,
worth $+\num{4.10}$~pp on CIFAR-10 and $+\num{3.18}$~pp on CIFAR-100 with the
optimiser held fixed. The
same step also carries essentially the whole cost increase,
\numrange{228}{272}~pp of \BP{}'s energy, while the other two steps together
move energy by at most \num{31}~pp.

\where{Section~5.3, Table~3, Figure~2; per-contrast intervals and
Holm-corrected $p$-values in \texttt{analysis/ladder\_analysis.json},
released with the paper.}

\subsection*{R2d: ``As backprop algorithms are more widely used, the
functionalities in existing software libraries could have been better
optimised \dots\ than forward-only algorithms.''}

\begin{comment}
As backprop algorithms are more widely used, the functionalities in existing
software libraries could have been better optimised for backprop than for
forward-only algorithms.
\end{comment}

\noindent
This is the objection we most wanted to be able to answer with data, and the
utilisation trace answers it. If kernel maturity decided the comparison, we
would expect \BP{} to be near saturation and the forward-only methods starved
of it. Neither holds. On the four ladder configurations, mean GPU utilisation
is \SI{11.5}{\percent} for \BP{} (maximum \SI{17.3}{\percent}) and
\SI{6.0}{\percent} for \MF{} (maximum \SI{8.0}{\percent}), and no method
exceeds \SI{20.3}{\percent}. Across the entire study, the wider \FF{} MLPs
included, no run exceeds \SI{38.0}{\percent} mean utilisation, still far
from saturation. Every arm is bound by kernel launch latency, not by
arithmetic throughput. This shows that \BP{} does not win by saturating the
A100 with more efficient arithmetic operations. It does not entirely exclude
an advantage from implementation maturity, and the revised paper says so in
Sections~5.5 and~6.

We also tested the specific optimisation a mature implementation would be
expected to have. \MF{}'s layer-sequential schedule recomputes the frozen
prefix at every stage, an $O(L^2)$ cost that an activation cache removes, so
we implemented the cache. It is bit-exact against the recomputing path:
maximum absolute difference in both loss and weights exactly \num{0} over
\num{1200} and \num{3600} optimiser steps. It bought no time
($+\SI{3.8}{\percent}$ and $+\SI{2.7}{\percent}$ wall-clock on device) and
cost \SI{374}{\percent} and \SI{404}{\percent} more per-process allocator memory,
which is \SI{242}{\mebibyte} and \SI{257}{\mebibyte} against
\SI{51}{\mebibyte}. Caching in host memory was worse on every axis. The
optimisation is available, and on this hardware it does not help.

Finally, the direction of the result does not need the objection resolved.
Under the matched data-pass protocol \MF{} is the \emph{cheaper} method, so any
library-maturity advantage \BP{} holds is working against the conclusion we
draw, not for it.

\where{Section~5.5, both paragraphs.}

% ============================================================
\section*{Reviewer 3}
% ============================================================

\subsection*{R3a: ``NVML \dots\ is never formally defined or explained.''}

\begin{comment}
NVML is used throughout but is never formally defined or explained.
\end{comment}

\noindent
Accepted. NVML is now expanded and defined on first use, in the introduction, as the
NVIDIA Management Library, the driver's own telemetry
interface, and the expansion is repeated in Sections~2 and~4. Section~4 has a dedicated \emph{Instruments} paragraph naming the
exact API call behind every quantity: \texttt{nvmlDeviceGetPowerUsage} for
power, from which energy is the trapezoidal integral of the \SI{5}{\hertz}
trace; \texttt{nvmlDeviceGetUtilizationRates} for utilisation;
\texttt{nvmlDeviceGetClockInfo} for SM clock;
\texttt{nvmlDeviceGetMemoryInfo} for device memory; and the host monotonic
clock for wall-clock time.

The same paragraph separates the two memory figures that the accepted version
conflated. The NVML reading is device-wide and includes a
${\approx}\SI{870}{\mebibyte}$ CUDA context that no algorithm controls. Every
memory number in the revised paper is instead \tmma{}, the per-process
allocator high-water mark, and we say so at each use. The driver version is
also stated, in ``Setup and statistics'', because it governs NVML
power-reading semantics.

\where{Sections~1 and~2; Section~4, ``Instruments'' and ``Setup and statistics''.}

\subsection*{R3b: ``The results for CaFo in Section~5.2 are presented purely
as text \dots\ should be summarised in an equivalent table.''}

\begin{comment}
The results for CaFo in Section 5.2 are presented purely as text, and should be
summarised in an equivalent table.
\end{comment}

\noindent
Accepted in substance, but we did not do exactly what was asked. To remain
within the \num{15}-page limit, we replaced the requested table with a compact
numerical summary in Section~5.6 and provide the complete table in the
repository. The revision is built around a single claim about \MF{},
and \FF{} and \CaFo{} are now context for it rather than co-equal subjects.
Giving each of them a results table would have cost roughly two pages of a
\num{15}-page limit for material that does not bear on the claim.

Instead, Section~5.6 states every \CaFo{} and \FF{} result numerically rather
than qualitatively: for \FF{}, the accuracy gap to \BP{} across its four arms
and its time, energy and memory as multiples of \BP{}'s; for \CaFo{}, the
accuracy gap on the MNIST-scale and CIFAR tasks and its energy as a multiple
of \BP{}'s; and the seed counts. \CaFo{} costs
\SIrange{3.0}{7.1}{\times} \BP{}'s energy, lands within \num{0.4}~pp of
\BP{} on the two MNIST-scale CNNs and loses \num{18.8} and \num{16.9}~pp on
CIFAR-10 and CIFAR-100; the direct-feedback-alignment variant narrows that to
\num{7.1} and \num{8.5}~pp at a further \SIrange{1.5}{2.6}{\times} the energy. The
complete per-configuration table for both algorithms is released as
\texttt{tables/ff\_cafo\_full.csv}, generated by
\texttt{analysis/ff\_cafo\_table.py} from \texttt{analysis/numbers.json}, which
the released analysis computes from the tidy table cited in the paper.

If the editors prefer the full table in the paper we will add it and cut
Section~5.6's prose to a cross-reference, which is close to page-neutral.

\where{Section~5.6; in the repository, \texttt{tables/\allowbreak ff\_cafo\_full.csv},
\texttt{analysis/\allowbreak numbers.json} and
\texttt{artifacts/\allowbreak tidy/\allowbreak runs.csv}.}

\subsection*{R3c: ``Because Mono-Forward places a cross-entropy loss at
every hidden layer, it inherently behaves like a deep supervision network
\dots\ add a standard backpropagation baseline that uses matching layer-wise
auxiliary losses. This is the only way to prove whether MF's gains stem from
its unique forward-only mechanism or simply from localised objective
constraints.''}

\begin{comment}
Because Mono-Forward places a cross-entropy loss at every hidden layer, it
inherently behaves like a deep supervision network. Add a standard
backpropagation baseline that uses matching layer-wise auxiliary losses. This
is the only way to prove whether MF's gains stem from its unique forward-only
mechanism or simply from localised objective constraints.
\end{comment}

\noindent
This is the most valuable comment we received and we have acted on it in full.
\BPDS{} is exactly the baseline described: \BP{} with \MF{}'s auxiliary
classifiers attached at every hidden layer and nothing detached, trained under
the same stopping rule, the same tuning budget and the same seeds.

The answer is decisive on the harder tasks and honest on the easier ones.
\MF{} beats \BPDS{} by \num{5.56}~pp $[\num{+5.37},\num{+5.75}]$ on CIFAR-10
and \num{6.10}~pp $[\num{+5.88},\num{+6.32}]$ on CIFAR-100, both over
\num{20} paired seeds with Holm-corrected $p<10^{-20}$. \BPDS{} is meanwhile
\emph{worse} than plain \BP{} on those two tasks, by \num{1.65} and
\num{1.06}~pp. Localised objectives therefore do not explain the gain; on the
tasks where \MF{} gains most, adding them to \BP{} costs accuracy. On the two
\num{2}$\times$\num{1000} MLPs the contrast is a tie, \MF{} leading \BPDS{} by
\num{0.11}~pp $[\num{-0.01},\num{+0.23}]$ on MNIST and trailing by
\num{0.06}~pp $[\num{-0.17},\num{+0.06}]$ on Fashion-MNIST, both EQUIVALENT
within the pre-registered margin. We report both halves in the paper: on the
small tasks the experiment does not distinguish an \MF{}-specific benefit from
auxiliary supervision, and on the harder ones auxiliary supervision does not
explain it.

Two further points on the strength of this evidence. The ladder separates the
readout from the detach as well, using a third rung, \MFJ{}, which has
\MF{}'s $\mathbf{M}_L$ readout and auxiliary losses but a global gradient. The step
from \MFJ{} to \MF{} detaches the gradient with the optimiser family held
fixed and is worth $+\num{4.10}$~pp on CIFAR-10 and $+\num{3.18}$~pp on
CIFAR-100. And the whole
analysis is pre-registered: the script fixing the paired test, the Holm
correction within metric family, the bootstrap intervals and the
$\pm\num{0.25}$~pp margin was committed before any ladder result was
inspected.

\where{Section~5.3, Table~3, Figure~2;
\texttt{scripts/analyze\_ablation\_ladder.py} for the pre-registration.}

\subsection*{R3d: the memory finding highlighted in the reviewer's summary,
and how to read the revision against it}

\begin{comment}
Crucially, the study debunks a common assumption in BP-free literature by
showing that the actual memory overhead introduced by algorithm-specific
components (such as local projection matrices) frequently counteracts or
entirely offsets the theoretical memory savings gained from omitting the
backward pass.
\end{comment}

\noindent
The reviewer singled this out, the revision now headlines a
\SIrange{24.6}{27.5}{\percent} memory \emph{saving} for \MF{}, and we want to
be explicit that this is a refinement of the finding rather than a retraction
of it.

The finding stands and the revision still carries it. \FF{} holds
\SIrange{0.91}{1.06}{\times} \BP{}'s allocator peak, so its dual positive and negative data
streams cancel most or all of the saving (Section~5.6). \MF{} at
\num{1.8}~M weights shows a \SI{0.3}{\percent} \emph{deficit}, which is
exactly the \SI{0.14}{\mebibyte} of extra projection parameters outweighing
what the absent backward pass saves (Section~5.2). What the revision adds is
the other end of the same axis: at \num{14.2}~M weights the activation saving
outgrows the projection overhead and the sign flips. The reviewer's point was
that the overhead is real and must be measured rather than assumed, and the
new Figure~1 is that point plotted, from a deficit at the small size to a
\SI{27.5}{\percent} saving at the large one.

One correction belongs here. The accepted version's \SI{5}{\percent} memory
figure came from the device-wide NVML counter, in which a
${\approx}\SI{870}{\mebibyte}$ CUDA context dilutes every difference. Every
memory number in the revision is \tmma{}, the per-process allocator
high-water mark. The instrument changed; the conclusion the reviewer
identified did not.

\where{Section~5.2 and Figure~1; Section~5.6; Section~4, ``Instruments''.}

% ============================================================
\section*{Reviewer 4}
% ============================================================

\subsection*{R4a: ``The key term `hardware-validated comparison' is unclear
and should be explained in more detail.''}

\begin{comment}
The key term ``hardware-validated comparison'' is unclear and should be
explained in more detail.
\end{comment}

\noindent
Accepted. The revision does not rely on the phrase to carry the meaning. The
introduction states what is measured and how, the \emph{Instruments} paragraph
of Section~4 names the API call behind each quantity, and Section~2 explains
why this matters rather than asserting that it does: analytical proxies such
as FLOP or parameter counts cannot see clock behaviour, utilisation or kernel
launch overhead, which is precisely where the differences in this study turn
out to live.

Section~5.5 makes the point concrete. A FLOP-based comparison of \MF{} and
\BP{} would predict a speed-up from the cache and see no difference in power.
The device says the opposite: the cache is bit-exact and saves no time, and the
entire energy result is a power difference at near-equal wall-clock time. That gap
between the proxy and the instrument is what ``hardware-validated'' is doing
in the title.

\where{Section~1; Section~2, ``Hardware and energy''; Section~4,
``Instruments''; Section~5.5.}

\subsection*{R4b: ``The statistical analysis is based on only three runs,
which is insufficient and can be considered only as preliminary evidence.''}

\begin{comment}
The statistical analysis is based on only three runs, which is insufficient and
can be considered only as preliminary evidence.
\end{comment}

\noindent
Accepted without reservation. Three seeds was insufficient and we have not
tried to defend it.

Every \MF{} contrast in the revised paper is over at least \num{20} seeds (the
\FF{} and \CaFo{} context runs use five to seven), and the
two most heavily used configurations carry \num{60} and \num{122}. The
matched data-pass protocol comparison of Table~2 uses \num{20} \MF{} seeds against
\num{60}, \num{122}, \num{20} and \num{20} \BP{} seeds in column order, paired by seed where
the arms share one. The ladder of Table~3 uses \numrange{20}{122} seeds per
rung. The study as a whole is almost \num{1000} runs.

The analysis itself changed as much as the seed count. The primary test is
paired by seed, because a seed fixes both the train and validation split and
the initialisation. A Welch unequal-variance test is computed alongside it in the released
analysis, never instead of it. Holm-Bonferroni correction is applied within each metric
family and never pooled across families. Confidence intervals come from
\num{10000} bootstrap resamples with a fixed seed. Most importantly, no parity
claim rests on a non-significant difference test: every one uses the two
one-sided tests (TOST) procedure, now cited and explained in Section~4, against
a $\pm\num{0.25}$~pp margin fixed in advance, and the
verdict vocabulary is restricted to DIFFERENT, EQUIVALENT and INCONCLUSIVE.

One scope statement belongs here. The ladder analysis is pre-registered, and
the paper says so. The matched data-pass protocol is \emph{not}
pre-registered; it is analysed with the same machinery and the same margins,
but the paper labels it a diagnostic rather than a confirmatory test.

\where{Section~4, ``Setup and statistics''; seed counts in the captions of
Tables~2 to~4.}

% ============================================================
\section*{Reviewer 5}
% ============================================================

\subsection*{R5a: ``The contribution is mainly evaluative rather than
algorithmic.''}

\begin{comment}
The contribution is mainly evaluative rather than algorithmic.
\end{comment}

\noindent
The observation is accurate and we have not tried to dress the paper up as
something else. What we have done is make the evaluative contribution carry a
mechanistic result rather than a league table.

The ablation ladder is the difference. It does not report that \MF{} differs
from \BP{}; it decomposes the difference into three single structural steps
and identifies the same one, the detach, as the largest contributor to both
the accuracy and the cost among the factors tested. That attribution is
falsifiable, it was pre-registered, and it contradicts the
reading a reader would most naturally have taken, namely that \MF{}'s
auxiliary losses do the work. It also produces a negative result of its own:
deep supervision on top of \BP{} is worse than \BP{} on both CIFAR tasks.

The second non-obvious contribution is methodological. The same algorithms support
opposite cost conclusions depending on the stopping rule, and the revision
shows both and quantifies the distance between them. We think the general
statement is worth making: an efficiency claim about a training algorithm is
uninterpretable without its stopping rule, because the rule is not a detail of
the measurement but a term in it. That is stated in Section~6 as a claim about
the field, not about this algorithm.

\where{Section~5.3; Section~6, ``Implications''.}

\subsection*{R5b: ``Small datasets, shallow architectures, and the use of
flattened CIFAR inputs for MLP experiments.''}

\begin{comment}
Small datasets, shallow architectures, and the use of flattened CIFAR inputs
for MLP experiments.
\end{comment}

\noindent
Accepted as a limitation, stated as one, and partly turned into a result.

On flattening: CIFAR inputs are flattened because that is the protocol of the
\MF{} source publication, and departing from it would make our numbers
incomparable with the ones we are testing. The important point is that
flattening is applied identically to \MF{} and to its \BP{} baseline, along
with the same augmentation, so it cannot favour either. Section~4 states
both: flattening follows the \MF{} protocol for the MLP experiments, and the
augmentation is applied identically to every algorithm and its baseline.

On scale: the revision states the bound rather than hoping the reader supplies
it. The headline percentages hold on the \num{3}$\times$\num{2000} MLPs, no
convolutional \MF{} was run, and both statements appear in Section~6. What we can offer instead of a larger model is the
direction of travel, measured: across a nearly eightfold range of trainable
weights every cost axis improves, and the memory advantage exists only at the
larger size, so the advantage is larger in the larger configurations tested,
with size confounded with depth and task. If the effect scaled the other way, that would be visible in
Figure~1, and it is not.

We would rather report a trend over the range we could afford to run
almost \num{1000} times than a single point at a larger scale with three seeds.

\where{Section~4; Section~5.2 and Figure~1; Section~6, limitations.}

\subsection*{R5c: ``Results are based on only three runs, and some of the
reported accuracy differences are modest.''}

\begin{comment}
Results are based on only three runs, and some of the reported accuracy
differences are modest.
\end{comment}

\noindent
The seed count is addressed under R4b: \numrange{20}{122} per \MF{} contrast,
almost \num{1000} runs in total.

The second half deserves its own answer, because the reviewer is right that
several differences are small, and the revision treats small differences as
such. Where a difference is modest we no longer call it a win. Every accuracy
parity claim is a two one-sided test against a $\pm\num{0.25}$~pp margin fixed
before the data were seen, and a contrast whose interval does not fit inside
the band is marked INCONCLUSIVE rather than being read as parity. That is why
the paper reports \MF{} as equivalent to \BP{} under the matched data-pass protocol on CIFAR-100
($+\num{0.02}$~pp, EQUIVALENT) but as \num{1.11}~pp \emph{behind} on CIFAR-10
(DIFFERENT), and does not average the two into a claim of parity.

The differences the paper does lean on are not modest. \MF{} over \BPDS{} is
\num{5.56} and \num{6.10}~pp; \MF{} over \BP{} at its own convergence is
\num{3.90} and \num{5.04}~pp; the memory saving is \num{82.2} and
\SI{74.4}{\mebibyte} with zero seed-to-seed variance, because the allocation
schedule of a fixed architecture at a fixed batch size is deterministic.

\where{Section~4, ``Setup and statistics''; Section~5.1; Section~5.3.}

\subsection*{R5d: ``A stronger comparison would include additional baselines,
such as deep-supervision or local-loss backpropagation variants, to better
isolate whether the gains come from the forward-only training mechanism or
from local auxiliary objectives.''}

\begin{comment}
A stronger comparison would include additional baselines, such as
deep-supervision or local-loss backpropagation variants, to better isolate
whether the gains come from the forward-only training mechanism or from local
auxiliary objectives.
\end{comment}

\noindent
Done, and the answer is the one reported under R3c and in the Summary above.
Two additional baselines were built, tuned and run at \numrange{20}{122} seeds
each:

\begin{description}[topsep=4pt,itemsep=2pt,parsep=0pt]
  \item[\BPDS{}] deep supervision: \BP{} with \MF{}'s layer-wise auxiliary
    classifiers, nothing detached. This is the local-auxiliary-objective
    control the reviewer asks for.
  \item[\MFJ{}] \MF{}'s learned $\mathbf{M}_L$ readout and auxiliary losses with a
    global gradient. This separates the extra parameters from the detach, so
    the readout cannot be credited with the locality effect.
\end{description}

\noindent
The separation is clear on the CIFAR tasks. Auxiliary objectives alone lose
\num{1.65} and \num{1.06}~pp against plain \BP{} on those two tasks. Adding the
readout recovers most or all of that, reaching $-\num{0.19}$ and $+\num{1.86}$~pp.
Only the final rung, which adds the detach, reaches $+\num{3.90}$ and
$+\num{5.04}$~pp. Among the factors tested, the gain is therefore attributable
mainly to gradient detachment rather than to the local objectives, which on
these two tasks lower accuracy when used alone.

\where{Section~5.3, Table~3, Figure~2.}

\subsection*{R5e: ``The energy measurements are useful, although specific to
one GPU platform.''}

\begin{comment}
The energy measurements are useful, although specific to one GPU platform.
\end{comment}

\noindent
Agreed, and the revision states the dependence in the specific form we think
it takes, rather than as a generic caveat.

All runs are on NVIDIA A100-SXM4-40GB with driver 595.71.05, and the driver
version is stated because it governs NVML power-reading semantics. The
limitation is not merely that another device might give another number; it is
that we can say which way. Board power at these utilisations is dominated by
the static and clock-domain components of a large A100 die, and the entire
energy result of Table~2 is a power difference at near-equal wall-clock time. A
smaller or better-saturated device would very likely narrow that power gap and
so narrow the energy saving. Section~6 says exactly this.

Two components of the result are more portable than the energy figure. The
memory saving is an allocator high-water mark and follows from the algorithm's
activation lifetime rather than from the device, and the pass multiplier that
dominates the own-convergence comparison is a property of the optimisation,
not of the hardware. We would expect both to survive a change of accelerator.

\where{Section~4, ``Setup and statistics''; Section~6, limitations.}

\subsection*{R5f: ``Clarify the limitations, and strengthen the discussion of
baselines and statistical robustness.''}

\begin{comment}
Clarify the limitations, and strengthen the discussion of baselines and
statistical robustness.
\end{comment}

\noindent
All three are addressed above and consolidated in Section~6, which now names
the binding limitation rather than listing caveats. That limitation is
convergence rate: \MF{} needs \SIrange{3.6}{4.4}{\times} \BP{}'s wall-clock
time to reach its own best accuracy, which converts a per-pass advantage of \SIrange{12.5}{19.9}{\percent} into an end-to-end penalty of
\SIrange{3.2}{4.0}{\times}. The paper shows that the two obvious remedies do
not close it and states it as the open problem, which is where we think work
on forward-only efficiency should go.

Three further limitations are bounded explicitly in Section~6: scope (the
headline percentages hold on the \num{3}$\times$\num{2000} MLPs, and there is
no memory advantage at \num{1.8}~M weights); accuracy under the matched data-pass protocol
(equivalent on one configuration of four, and \num{1.11}~pp behind on
CIFAR-10, so the matched data-pass protocol result is a cost result and not an accuracy
result); and hardware (R5e). Pre-registration (the ladder is, the matched
data-pass protocol is not) is stated in Section~4. Baselines are R2c, R3c and R5d. Statistical robustness is
R4b and R5c.

\where{Section~6, ``The open problem: convergence rate, not per-unit cost''.}

\bigskip
\noindent
We are grateful for the review. R3c and R5d in particular produced the
experiment that this version is built on, and it is a better paper than the one
that was accepted.

\end{document}
